\chapter{Descripción general del sistema}

\section{Función general}

La Raspberry Pi constituye el controlador principal de SmartGate. En ella se ejecuta una aplicación desarrollada en Python que integra la cámara, los modelos, la lógica de autorización, las entradas y salidas GPIO y la comunicación con Supabase. Esta elección permite realizar localmente las decisiones críticas: la detección, la validación de vida, el reconocimiento facial y la apertura del portón no dependen de que exista una conexión permanente a Internet.

\section{Componentes del sistema y su interacción}

El sistema se compone de cuatro bloques: la interfaz web de administración, la base de datos en Supabase, la Raspberry Pi que realiza el reconocimiento y la interfaz electrónica que actúa sobre el portón. La interfaz web y la Raspberry Pi trabajan sobre la misma fuente de datos: cuando un usuario es creado o actualizado desde la web, la modificación queda almacenada en Supabase y puede ser consultada por la Raspberry Pi. En sentido inverso, los eventos enviados por el dispositivo quedan disponibles en la interfaz para su supervisión. La figura~\ref{fig:interaccion} resume esta interacción.

\begin{figure}[H]
    \centering
    \begin{tikzpicture}[node distance=1.5cm]
        \node[caja] (web) {\textbf{Interfaz web}\\ Administración de usuarios\\ Carga de imágenes y consultas};
        \node[caja, below=of web] (supa) {\textbf{Supabase}\\ Base de datos PostgreSQL\\ Almacenamiento y eventos};
        \node[caja, below=of supa] (rpi) {\textbf{Raspberry Pi}\\ Reconocimiento facial\\ Procesamiento y decisión};
        \node[caja, below=of rpi] (sg) {\textbf{SmartGate}\\ Interfaz electrónica\\ Control de acceso físico};
        \draw[dflecha] (web) -- node[etiqueta]{API y consultas} (supa);
        \draw[dflecha] (supa) -- node[etiqueta]{API HTTPS: usuarios y registros} (rpi);
        \draw[dflecha] (rpi) -- node[etiqueta]{Orden de apertura y estado del portón} (sg);
    \end{tikzpicture}
    \caption{Interacción entre la interfaz web, Supabase, la Raspberry Pi y el sistema de acceso.}
    \label{fig:interaccion}
\end{figure}

Cada uno de estos bloques se desarrolla en los capítulos siguientes: la interfaz con el portón en el capítulo~\ref{cap:interfaz}, los modelos de reconocimiento en el capítulo~\ref{cap:modelos}, el programa de la Raspberry Pi en el capítulo~\ref{cap:programa} y la base de datos, la interfaz web y el funcionamiento sin conexión en el capítulo~\ref{cap:datos}.

\section{Organización del desarrollo}

El desarrollo de SmartGate se realizó de forma escalonada. Desde el inicio se separó el trabajo en dos niveles: un producto mínimo viable (MVP), encargado de demostrar el funcionamiento principal del sistema, y un conjunto de funcionalidades adicionales que solamente se abordarían una vez que el núcleo estuviera integrado y probado. Esta separación permitió mantener un objetivo funcional durante todo el proyecto y evitó que las mejoras opcionales comprometieran la entrega principal.

El trabajo comenzó con la investigación de modelos de detección facial, reconocimiento y antispoofing, junto con el análisis de la Raspberry Pi, la cámara y la electrónica necesaria para interactuar con el portón. A partir de esa investigación se definió una arquitectura modular y se utilizaron modelos previamente entrenados para construir una primera versión del pipeline. Luego se incorporó el control físico y, una vez validado el ciclo completo, se avanzó sobre la base de datos, la interfaz web, la persistencia offline y las herramientas relacionadas con el entrenamiento de modelos propios.

\begin{figure}[H]
    \centering
    \begin{tikzpicture}[
        phase/.style={rectangle, rounded corners, draw=black, thick, align=center, text width=2.55cm, minimum height=1.45cm, font=\small, fill=blue!8},
        extra/.style={rectangle, rounded corners, draw=black, thick, align=center, text width=2.55cm, minimum height=1.45cm, font=\small, fill=green!10},
        pending/.style={rectangle, rounded corners, draw=black, thick, align=center, text width=2.55cm, minimum height=1.45cm, font=\small, fill=orange!15}
    ]
        \node[phase] (research) at (0,0) {Investigación\\ y definición};
        \node[phase] (mvp) at (3.8,0) {Implementación\\ del MVP};
        \node[extra] (extras) at (7.6,0) {Extras\\ operativos};
        \node[pending] (training) at (11.4,0) {Dataset y\\ entrenamiento};

        \draw[flecha] (research) -- (mvp);
        \draw[flecha] (mvp) -- (extras);
        \draw[flecha] (extras) -- (training);
    \end{tikzpicture}
    \caption{Desarrollo escalonado de SmartGate. La última etapa se inició, pero no pudo completarse.}
    \label{fig:desarrollo-escalonado}
\end{figure}

La figura~\ref{fig:desarrollo-escalonado} resume las etapas: investigación y definición de la arquitectura, el hardware y los modelos; implementación del MVP, que comprende el pipeline y el control físico; extras operativos, como Supabase, la interfaz web y la persistencia; y finalmente el dataset y el entrenamiento propio, etapa que se inició pero no se completó. El detalle de lo implementado en cada nivel se presenta en el capítulo~\ref{cap:desarrollo}.
